%==============================================================================
% فصل پنجم: بحث و نتیجه‌گیری
% بحث، نتیجه‌گیری، محدودیت‌ها، پیشنهادها
%==============================================================================
\ChapterStart
\chapter{بحث و نتیجه‌گیری}
\label{ch:discussion}

\section{بحث}
\label{sec:discussion}


پژوهش حاضر با هدف تعیین میزان پایبندی به اقدامات غیردارویی پیشگیری ثانویه و شناسایی عوامل مرتبط با آن در ۱۵۰ بیمار مبتلا به انفارکتوس میوکارد مراجعه‌کننده به درمانگاه قلب بیمارستان شهید مدنی خرم‌آباد انجام شد. مهم‌ترین یافته مطالعه آن بود که با وجود پایبندی قابل قبول بخشی از بیماران به مؤلفه‌های مورد بررسی، تنها ۲۶ نفر ($17.3$ درصد) به‌طور هم‌زمان به همه اقدامات غیردارویی مورد بررسی پایبند بودند. این نتیجه نشان می‌دهد که مسئله اصلی در مراقبت پس از \lr{MI}، فقدان رفتارهای سالم نیست، بلکه دشواری حفظ هم‌زمان مجموعه‌ای از رفتارهای توصیه‌شده است. به بیان دیگر، موفقیت در یک حیطه لزوما به معنای تحقق کامل پیشگیری ثانویه نیست و ارزیابی تک‌بعدی ممکن است وضعیت واقعی بیمار را متفاوت از آنچه هست نشان دهد.

این یافته با الگوی غالب شواهد علمی بین‌المللی هم‌سو است. در مطالعه آینده‌نگر مورنو\footnote{\lr{Moreno}} و همکاران، $84.6$ درصد بیماران تا ۱۲ ماه پس از \lr{MI} دست‌کم به یکی از توصیه های دارویی، فعالیت بدنی یا رژیم غذایی پایبند نبودند\cite{Moreno2024NonAdherence}. همچنین، در مطالعه کاتسوا و همکاران، تنها ۲۹ درصد بیماران به هر سه توصیه اصلی سبک زندگی (ترک سیگار، کاهش وزن و فعالیت بدنی) دست یافته بودند\cite{DeBacquer2022Lifestyle,Kotseva2019EUROASPIRE}؛ و در یک مطالعه همگروهی شش‌ساله در فرانسه، فقط ۱۰ درصد بیماران همه توصیه های درمانی، کنترل عوامل خطر و سبک زندگی را محقق کرده بودند\cite{DibaoDina2018Adherence}. تفاوت میان برآوردها می‌تواند ناشی از تفاوت در تعریف پایبندی کامل، زمان سپری‌شده از \lr{MI}، ویژگی‌های جمعیتی، دسترسی به خدمات بازتوانی و ابزارهای سنجش باشد؛ بااین‌حال، جهت کلی یافته‌ها یکسان است و نشان دهنده وجود شکافی قابل توجه میان توصیه بالینی و رفتار روزمره بیماران می باشد.

در مطالعه حاضر، ۲۷ بیمار (۱۸ درصد) مصرف‌کننده دخانیات بودند. این شاخص در مقایسه با سایر اقدامات پیشگیری ثانویه وضعیت نسبتا مطلوب‌تری داشت. افزون بر این، خودگزارش‌دهی مصرف دخانیات می تواند باعث گزارش کمتر مصرف شود. مشاهده وضعیت بهتر در حیطه دخانیات نسبت به رفتارهای پیچیده‌تری مانند فعالیت بدنی یا پایبندی هم‌زمان، با شواهد موجود نیز سازگار است؛ اغلب مطالعات بین‌المللی نشان داده‌اند که ترک دخانیات در جمعیت‌های مورد مطالعه بیش از ورزش منظم یا مشارکت در بازتوانی تحقق می‌یابد\cite{DeBacquer2022Lifestyle,Urbinati2015BLITZ4}.

در تحلیل چندمتغیره، احتمال مصرف دخانیات در مردان و افراد بیکار، بازنشسته یا خانه‌دار بیشتر بود. از آنجا که هیچ‌یک از زنان حاضر در جمعیت مورد بررسی دخانیات مصرف نمی‌کردند، این نتیجه باید با احتیاط تفسیر شود. ارتباط اشتغال با مصرف دخانیات می‌تواند از مواجهه‌های محیط کار، شبکه‌های اجتماعی، فشار شغلی یا الگوهای جنسیتی اشتغال تأثیر پذیرفته باشد، ازاین‌رو، یافته حاضر بیش از آنکه یک رابطه علّی را اثبات کند، گروهی را برای غربالگری و حمایت هدفمند مشخص می‌سازد.

میانه فعالیت بدنی بیماران ۲۴۰ دقیقه در هفته با دامنه میان‌چارکی ۱۲۰ تا ۴۵۰ دقیقه بود. اگرچه این مقدار در سطح گروهی می‌تواند مطلوب به نظر برسد، پراکندگی زیاد آن نشان می‌دهد که تجربه بیماران یکسان نبوده و 46 نفر ($30.7$ درصد) به حد توصیه‌شده ۱۵۰ دقیقه فعالیت بدنی با شدت متوسط در هفته نرسیده‌اند. افزون بر این، مقدار زمان گزارش‌شده به‌تنهایی اطلاعات کاملی درباره شدت، تداوم، ایمنی و پایداری فعالیت در طول زمان ارائه نمی‌کند. شواهد بین المللی نیز فعالیت بدنی را یکی از چالشی‌ترین حیطه‌های پیشگیری ثانویه معرفی کرده‌اند. در مطالعه مورنو و همکاران، $28.6$ درصد بیماران در طول ۱۲ ماه به هدف فعالیت بدنی پایبند نبودند\cite{Moreno2024NonAdherence}. در یک مطالعه آلمانی فقط حدود یک‌سوم بیماران ورزش منظم داشتند\cite{Zeymer2024PatientKnowledge} و در یک مطالعه چند‌مرکزی اروپایی نیز تنها ۴۰ درصد به هدف فعالیت بدنی دست یافتند\cite{DeBacquer2022Lifestyle}. تفاوت نرخ‌ها با مطالعه حاضر می‌تواند به ابزار سنجش، تعریف فعالیت بدنی، زمان پیگیری و تفاوت‌های محیطی و فرهنگی مربوط باشد.

در مدل چندمتغیره، مردان با احتمال بیشتری به فعالیت بدنی پایبند بودند. این یافته با گزارش برخی مطالعات درباره مشارکت کمتر زنان در فعالیت بدنی و بازتوانی قلبی هم‌خوانی دارد\cite{Moreno2024NonAdherence,DeBacquer2022Lifestyle,DibaoDina2018Adherence} و می‌تواند با مسئولیت‌های خانوادگی، فرصت کمتر برای ورزش، نگرانی از ایمنی فعالیت، دسترسی محدودتر به فضای مناسب و حمایت اجتماعی کمتر مرتبط باشد.

وضعیت اشتغال نیز با فعالیت بدنی ارتباط داشت و افراد بیکار، بازنشسته یا خانه‌دار پایبندی کمتری داشتند. ممکن است اشتغال با تحرک روزانه، ساختار منظم‌تر زندگی یا وضعیت عملکردی بهتر همراه باشد؛ از سوی دیگر، «اثر کارگر سالم»\footnote{\lr{Healthy worker effect}} نیز محتمل است، یعنی افرادی که توان جسمی بیشتری دارند امکان ادامه اشتغال بیشتری پیدا می‌کنند. بنابراین، نمی‌توان اشتغال را به‌تنهایی علت فعالیت بدنی بیشتر دانست.

نوع درمان نیز با فعالیت بدنی ارتباط مستقل داشت. بیمارانی که درمان دارویی دریافت کرده بودند، در مقایسه با بیماران تحت آنژیوپلاستی، با احتمال کمتری به فعالیت بدنی پایبند بودند، درحالی‌که تفاوت جراحی بای‌پس عروق کرونر با آنژیوپلاستی معنادار نبود. این ارتباط ممکن است تفاوت در شدت بیماری، وضعیت عملکردی، احساس بهبودی کاذب بیمار پس از بازگشایی عروق، نوع توصیه‌های هنگام ترخیص یا شدت پیگیری را منعکس کند. از آنجا که تخصیص نوع درمان در این مطالعه تصادفی نبوده است، احتمال مخدوش‌شدن نتایج بر اثر شدت بیماری و سایر عوامل بالینی وجود دارد.

در تحلیل دومتغیره، سواد سلامت، تعداد بیماری‌های همراه و پایبندی دارویی با فعالیت بدنی ارتباط داشتند، اما این روابط پس از ورود هم‌زمان متغیرها به مدل چندمتغیره معنی‌دار نبودند. این تغییر نشان می‌دهد که بخشی از ارتباط خام این متغیرها احتمالا با تفاوت‌های جمعیت‌شناختی یا درمانی توضیح داده می‌شود.

از نظر تغذیه، ۹۰ بیمار (۶۰ درصد) بر اساس پرسشنامه \lr{MEDAS} به رژیم غذایی مدیترانه‌ای پایبند بودند. این میزان در محدوده مطالعات پیشین قرار می‌گیرد؛ در مطالعه مورنو و همکاران، $43.2$ درصد بیماران در طول ۱۲ ماه به رژیم مدیترانه‌ای پایبند نبودند\cite{Moreno2024NonAdherence}. هم‌سویی نسبی این دو برآورد، با وجود تفاوت محیط فرهنگی و الگوی غذایی، نشان می‌دهد حفظ یک الگوی غذایی سالم پس از رویداد حاد برای بخش قابل‌توجهی از بیماران دشوار است. رژیم غذایی برخلاف یک توصیه منفرد، مجموعه‌ای از انتخاب‌های تکرارشونده روزانه است و از قیمت و دسترسی مواد غذایی، عادات خانوادگی، دانش تغذیه‌ای و امکان آماده‌سازی غذا تأثیر می‌پذیرد\cite{Krezel2024Nutrition,Deslippe2023Barriers,Katsi2025Diet}.

در تحلیل‌های دومتغیره، تحصیلات، محل سکونت، سواد سلامت و پایبندی دارویی با پایبندی غذایی ارتباط داشتند؛ اما در مدل چندمتغیره، تنها وضعیت اقتصادی ارتباط مستقل خود را حفظ کرد به نحوی که وضعیت اقتصادی بهتر با افزایش احتمال پایبندی به رژیم غذایی توصیه شده همراه بود. این یافته از نقش مسائل اقتصادی در تبدیل دانش به رفتار حمایت می‌کند. ممکن است بیمار توصیه‌های تغذیه‌ای را بداند، اما توان خرید مستمر میوه، سبزی، ماهی، مغزها یا روغن‌های مناسب را نداشته باشد.

مطالعات انجام‌شده در جمعیت‌های کم‌برخوردار نیز وضعیت اقتصادی نامطلوب و دشواری دسترسی را از موانع پایبندی به پیشگیری ثانویه گزارش کرده‌اند\cite{Mohandas2025RecurrentMI}. عدم وجود ارتباط معنی‌دار بین سواد سلامت و پایبندی به توصیه های تغذیه ای در مدل نهایی به معنای بی‌اهمیت‌بودن آن‌ها نیست؛ بلکه ممکن است تاثیر آن از مسیر تحصیلات، وضعیت اقتصادی یا پایبندی عمومی به درمان اعمال شود.

در این مطالعه بر اساس پرسشنامه \lr{PHQ-4}، ۶۹ بیمار (۴۶ درصد) وضعیت روان‌شناختی طبیعی داشتند. بنابراین، بیش از نیمی از نمونه‌ها درجاتی از دیسترس روان‌شناختی را گزارش کردند. این یافته از بار قابل‌توجه مشکلات روانشناختی پس از \lr{MI} حکایت دارد. مطالعات اخیر نیز افسردگی، اضطراب، ترس از بیماری مجدد، کاهش استقلال فردی و تصور بیمار از بیماری را از عوامل مهم در تداوم رفتارهای سالم پس از \lr{MI} دانسته‌اند\cite{AlFaraj2025MentalHealth}. در یک مطالعه مشابه، نشانه‌های افسردگی و بار مراقبتی با عدم پایبندی به مجموعه اقدامات پیشگیرانه ارتباط داشت. همچنین در یک مطالعه طولی، حمایت اجتماعی درک‌شده با پایبندی بلندمدت بهتر همراه بود\cite{Moreno2024NonAdherence,Nachshol2020Psychosocial}. بر این اساس، سلامت روان را نباید پیامدی فرعی و جدا از پیشگیری قلبی تلقی کرد؛ وضعیت روان‌شناختی می‌تواند بر انگیزش، برنامه‌ریزی، خودمراقبتی و توان حفظ تغییرات رفتاری اثر بگذارد.

در مدل تعدیل‌شده سلامت روان‌شناختی، تأهل با وضعیت بهتر و جنس مؤنث با وضعیت نامطلوب‌تر همراه بود. نقش محافظتی تأهل می‌تواند با حمایت عاطفی و عملی همسر، کمک در مصرف دارو و تغذیه، همراهی در مراجعه و کاهش انزوای اجتماعی توضیح داده شود؛ هرچند کیفیت رابطه و حمایت اجتماعی در مطالعه حاضر مستقیما سنجیده نشده است. وضعیت نامطلوب‌تر زنان نیز با شواهدی که بار بیشتر اضطراب و افسردگی پس از حوادث قلبی را در زنان گزارش کرده‌اند هم‌خوان است و می‌تواند تحت تأثیر تفاوت‌های اجتماعی، اقتصادی، مسئولیت‌های مراقبتی و دسترسی به حمایت قرار گیرد\cite{Flagtvedt2025Gender,Sullivan2022GenderDisparities,Serpytis2018Gender}. همچنین، سطح ادراک بیماری پایین‌تر به‌طور مستقل با سلامت روان‌شناختی ضعیف‌تر همراه بود. این یافته با مدل خودتنظیمی لونتال سازگار است؛ بر اساس این مدل، بازنمایی شناختی و هیجانی بیمار از ماهیت، پیامدها، مدت و کنترل‌پذیری بیماری بر واکنش عاطفی و راهبردهای مقابله‌ای او اثر می‌گذارد. بااین‌حال، در این مطالعه به دلیل هم‌زمانی سنجش ادراک بیماری و سلامت روان، جهت رابطه مشخص نیست.


\begin{figure}[H]
  \centering
  \begin{tikzpicture}[
    x=1.75cm,
    y=1.7cm,
    font=\fontsize{11pt}{14pt}\selectfont,
    box/.style={
      draw=black,
      line width=0.7pt,
      fill=white,
      minimum height=1.05cm,
      inner xsep=5pt,
      inner ysep=3pt,
      align=center
    },
    perception/.style={box,minimum width=3.45cm},
    coping/.style={box,minimum width=2.35cm},
    outcome/.style={box,minimum width=2.55cm},
    disease/.style={box,minimum width=2.65cm},
    flow/.style={line width=0.85pt,-{Latex[length=2.2mm,width=1.7mm]}}
  ]
    \node[disease] (disease) at (0,0)
      {\shortstack{\TikzRTL{ویژگی‌های}\\\TikzRTL{بیماری}}};

    \node[perception] (cognitive) at (2.55,1.25)
      {\shortstack{\TikzRTL{ادراک‌های شناختی}\\\TikzRTL{بیماری}}};
    \node[perception] (emotional-perception) at (2.55,-1.25)
      {\shortstack{\TikzRTL{ادراک‌های هیجانی}\\\TikzRTL{بیماری}}};

    \node[coping] (cognitive-coping) at (5.25,1.25)
      {\shortstack{\TikzRTL{راهبردهای}\\\TikzRTL{مقابله‌ای}}};
    \node[coping] (emotional-coping) at (5.25,-1.25)
      {\shortstack{\TikzRTL{راهبردهای}\\\TikzRTL{مقابله‌ای}}};

    \node[outcome] (illness-outcome) at (7.75,1.25)
      {\shortstack{\TikzRTL{پیامدهای}\\\TikzRTL{بیماری}}};
    \node[outcome] (emotional-outcome) at (7.75,-1.25)
      {\shortstack{\TikzRTL{پیامدهای}\\\TikzRTL{هیجانی}}};

    % Two principal causal paths.
    \draw[flow] (disease.north east) -- (cognitive.west);
    \draw[flow] (disease.south east) -- (emotional-perception.west);
    \draw[flow] (cognitive.east) -- (cognitive-coping.west);
    \draw[flow] (cognitive-coping.east) -- (illness-outcome.west);
    \draw[flow] (emotional-perception.east) -- (emotional-coping.west);
    \draw[flow] (emotional-coping.east) -- (emotional-outcome.west);

    % Reciprocal relation between illness and emotional outcomes.
    \draw[
      line width=0.85pt,
      {Latex[length=2.2mm,width=1.7mm]}-{Latex[length=2.2mm,width=1.7mm]}
    ] (illness-outcome.south) -- (emotional-outcome.north);

    % Upper feedback route from illness outcomes.
    \draw[line width=0.85pt]
      (illness-outcome.north) -- (7.75,2.10) -- (0,2.10);
    \draw[flow] (0,2.10) -- (disease.north);
    \draw[flow] (2.55,2.10) -- (cognitive.north);
    \draw[flow] (5.25,2.10) -- (cognitive-coping.north);

    % Lower feedback route from emotional outcomes.
    \draw[line width=0.85pt]
      (emotional-outcome.south) -- (7.75,-2.10) -- (0,-2.10);
    \draw[flow] (0,-2.10) -- (disease.south);
    \draw[flow] (2.55,-2.10) -- (emotional-perception.south);
    \draw[flow] (5.25,-2.10) -- (emotional-coping.south);
  \end{tikzpicture}
  \caption{مدل خودتنظیمی لونتال}
  \label{fig:illness-perception-model-fa}
\end{figure}

یکی از یافته‌های مهم مطالعه، ارتباط سواد سلامت با چند پیامد در تحلیل‌های دومتغیره بود. بیماران دارای سواد سلامت کافی با احتمال بیشتری، فعالیت بدنی کافی، رژیم غذایی مناسب و پایبندی کلی بالاتری داشتند و سواد سلامت ناکافی با دیسترس روان‌شناختی بیشتر همراه بود. این الگو با مطالعه لو و همکاران در بیماران مبتلا به بیماری عروق کرونر هم‌خوانی دارد که در آن سواد سلامت محدود، استقلال فردی پایین و اطلاعات ناکافی بیماری با پایبندی کمتر به رفتارهای سالم مرتبط بود\cite{Lu2020HealthLiteracy}. همچنین، مطالعه زیمر\footnote{\lr{Zeymer}} و همکاران نشان داد که احساس آگاهی بیماران لزوما با دانش عینی همراه نیست\cite{Zeymer2024PatientKnowledge}. با وجود این، سواد سلامت پس از تعدیل در مدل‌های فعالیت بدنی، رژیم غذایی، سلامت روان و پایبندی کلی ارتباط مستقلی نشان نداد. این نتیجه می‌تواند ناشی از هم‌پوشانی سواد سلامت با تحصیلات و وضعیت اقتصادی، اندازه محدود نمونه یا نقش میانجی آن باشد. بنابراین، یافته نهایی را نباید لزوما به‌صورت رد نقش سواد سلامت تفسیر کرد؛ بلکه به نظر می‌رسد سواد سلامت در شبکه‌ای از عوامل فردی و اجتماعی عمل می‌کند.

در مدل نهایی پایبندی کامل، تنها سطح تحصیلات ارتباط مستقل و معنادار داشت؛ بیماران دارای تحصیلات دانشگاهی در مقایسه با افراد دارای تحصیلات دیپلم و پایین‌تر شانس بیشتری برای پایبندی به همه اقدامات پیشگیرانه داشتند. این یافته با گزارش مطالعاتی که آموزش، دانش بیماری و سواد سلامت را با رفتارهای سالم پس از \lr{MI} مرتبط دانسته‌اند هم‌سو است\cite{Lu2020HealthLiteracy,Jaim2025APCU}. تحصیلات می‌تواند فهم توصیه‌ها، جست‌وجوی اطلاعات، ارتباط با کادر درمان، حل مسئله و دسترسی به منابع را تسهیل کند. بااین‌حال، تحصیلات احتمالا شاخصی برای مجموعه‌ای از مزیت‌های اجتماعی و شناختی است و خود به‌تنهایی سازوکار تغییر رفتار محسوب نمی‌شود. نسبت شانس تعدیل‌شده $6.574$ (فاصله اطمینان $1.93$ تا $22.29$)، باوجود معناداری آماری، عدم‌قطعیت قابل‌توجهی دارد؛ ازاین‌رو، اندازه دقیق اثر باید در مطالعاتی با حجم نمونه‌ بزرگ‌تر تأیید شود.

برخلاف تحلیل‌های دومتغیره، جنسیت، تأهل، سواد سلامت، ادراک بیماری و پایبندی دارویی در مدل نهایی پایبندی کلی معنادار نبودند. این تفاوت میان نتایج اولیه و تعدیل‌شده اهمیت تعامل و هم‌پوشانی عوامل را نشان می‌دهد. ممکن است یک عامل برای یک رفتار مشخص اهمیت داشته باشد، اما در شاخص ترکیبی سخت‌گیرانه‌ای که مستلزم موفقیت هم‌زمان در همه حیطه‌هاست، اثر مستقل آن آشکار نشود. افزون بر این، تعداد کم بیماران کاملا پایبند توان آماری مدل را محدود می‌کند.

هیچ‌یک از ۱۵۰ بیمار سابقه شرکت در برنامه های بازتوانی قلبی را نداشتند. این یافته شاید مهم‌ترین چالش نظام ارائه خدمت در مطالعه حاضر باشد. مرور شواهد نشان می‌دهد بازتوانی قلبی نه‌تنها یک مداخله منفرد، بلکه بستری برای آموزش، ورزش ایمن، حمایت روان‌شناختی، اصلاح تغذیه و ترک دخانیات است. در مطالعه مورنو و همکاران فقط ۴۵ درصد بیماران ارجاع‌شده برنامه را آغاز و $33.7$ درصد آن را کامل کردند؛ در یک مطالعه در اتریش میزان پایبندی به برنامه بازتوانی $13.4$ درصد بود\cite{Moreno2024NonAdherence,Alhabib2026Saudi,Hammer2024Adherence}. مطالعات \lr{BLITZ-4} و پیگیری شش‌ساله فرانسه نیز برنامه بازتوانی را با دستیابی بهتر به اهداف سبک زندگی مرتبط دانسته‌اند\cite{DibaoDina2018Adherence,Urbinati2015BLITZ4}.

مداخلات مؤثر باید علاوه بر آموزش، محدودیت‌های اقتصادی، روان‌شناختی، خانوادگی و ساختاری را نیز در نظر گیرند. شواهد مداخله‌ای نیز نشان داده‌اند که آموزش بیمار می‌تواند رفتار‌‌‌های ترک دخانیات، فعالیت بدنی کافی، رژیم غذایی مناسب و پایبندی درمانی را بهبود دهد، به‌ویژه هنگامی که پیگیری آن تداوم داشته باشد\cite{Shi2023Education,Ivers2020ISLAND}.



\section{پیشنهادها برای پژوهش‌های آینده}

\begin{itemize}
\item مطالعات آینده به‌صورت چندمرکزی و با نمونه بزرگ‌تر انجام شوند تا پایداری عوامل مرتبط، به‌ویژه برآوردهای مربوط به جنسیت، اشتغال، تحصیلات و نوع درمان، با دقت بیشتری سنجیده و تعمیم‌پذیری نتایج افزایش یابد.

\item طراحی طولی از زمان بستری تا چند مقطع پس از ترخیص به‌کار رود تا تغییر رفتارها، زمان افت پایبندی و تقدم زمانی عوامل مرتبط مشخص شود.

\item برای زنان، موانع فعالیت بدنی به‌طور اختصاصی ارزیابی شود و برنامه‌های ایمن، تدریجی، متناسب با مسئولیت‌های خانوادگی و قابل‌اجرا در منزل یا محیط‌های در دسترس طراحی گردد. آموزش درباره شدت مجاز فعالیت و علائم هشدار می‌تواند ترس از حرکت پس از \lr{MI} را کاهش دهد.
\end{itemize}

\section{نتیجه‌گیری}

نتایج این مطالعه نشان داد که عوامل مرتبط با پایبندی در حیطه‌های مختلف یکسان نبودند. بر اساس یافته‌ها، پایبندی غیردارویی پدیده‌ای چندبعدی است و از برهم‌کنش توانمندی‌های فردی، شرایط اجتماعی‌اقتصادی، وضعیت روانی و سازمان خدمات سلامت تأثیر می‌پذیرد.

به همین دلیل، پاسخ مناسب به عدم پایبندی نباید به توصیه های شفاهی و یکسان برای همه بیماران محدود شود. ارزیابی منظم و حیطه‌به‌حیطه رفتارها، شناسایی بیماران دارای تحصیلات پایین‌تر یا محدودیت اقتصادی، توجه ویژه به فعالیت بدنی زنان، غربالگری سلامت روان و ادراک بیماری، و ایجاد مسیر قابل‌دسترس بازتوانی قلبی می‌تواند پایه یک برنامه پیشگیری ثانویه جامع‌تر باشد. با توجه به طرح مقطعی پژوهش، این نتایج روابط آماری را نشان می‌دهند و برای استنباط علی یا تعمیم به همه بیماران مبتلا به \lr{MI} باید با احتیاط به‌کار روند.
